% Section 2: Methodology

\section{Methodology}

\subsection{xRFM: Algorithm and Theoretical Background}

xRFM is a tabular feature-learning model based on the Average Gradient Outer Product (AGOP)~\citep{beaglehole2025xrfm,radhakrishnan2024mechanism}. For a learned prediction function $f:\mathbb{R}^{d}\rightarrow\mathbb{R}$, the AGOP is defined as
\begin{equation}
    G_f=\mathbb{E}_{x}\!\left[\nabla f(x)\nabla f(x)^{\top}\right],
    \label{eq:agop}
\end{equation}
where $\nabla f(x)$ measures the sensitivity of the prediction to each input feature. Large diagonal values of $G_f$ indicate features that strongly affect the learned function. In this project, we use the diagonal AGOP setting because it is cheaper than storing the full $d\times d$ matrix and gives direct feature-level importance scores.

The key intuition is that AGOP provides a supervised view of feature relevance. PCA identifies high-variance input directions without using the target variable~\citep{pearson1901pca}; mutual information measures feature-target dependence mostly at the individual-feature level~\citep{kraskov2004estimating}; and permutation importance measures post-hoc test-performance degradation~\citep{fisher2019all}. By contrast, AGOP is derived from the gradients of the fitted model, so it reflects the directions that the prediction function actually uses. xRFM iteratively fits a base model, estimates AGOP, reweights the feature space, and refits the model to focus learning on task-relevant directions.

\subsection{Experimental Design}

\subsubsection{Datasets}

We evaluate xRFM and the baseline models on five public tabular datasets covering regression and binary classification tasks. These datasets vary in sample size, dimensionality, feature type, and class imbalance, allowing us to examine predictive performance, scalability, and interpretability under different tabular conditions~\citep{hamidieh2018data,sakar2019real,fedesoriano2021stroke}. Table~\ref{tab:datasets} summarises the final data statistics after preprocessing.

\begin{table}[H]
  \centering
  \caption{Dataset statistics. $n_\text{train}$, $n_\text{val}$, and $n_\text{test}$ denote the sizes of the training, validation, and test splits respectively; $d$ is the number of features after preprocessing.}
  \label{tab:datasets}
  \begin{tabular}{lllrrrr}
    \toprule
    Dataset & Task & Target & $n_\text{train}$ & $n_\text{val}$ & $n_\text{test}$ & $d$ \\
    \midrule
    Diamonds            & Regression     & price          & 32{,}365 & 10{,}789 & 10{,}789 & 26 \\
    Superconductivity   & Regression     & critical\_temp & 12{,}757 &  4{,}253 &  4{,}253 & 81 \\
    Online Shoppers     & Classification & Revenue        &  7{,}398 &  2{,}466 &  2{,}466 & 27 \\
    Stroke              & Classification & stroke         &  3{,}066 &  1{,}022 &  1{,}022 & 21 \\
    HR Attrition        & Classification & Attrition      &    882   &    294   &    294   & 51 \\
    \bottomrule
  \end{tabular}
\end{table}

\subsubsection{Preprocessing and Splitting}

All experiments use a fixed random seed of 42 and a 60/20/20 train-validation-test split. Missing numerical values are imputed using the median, while missing categorical values are replaced with a constant \texttt{Missing} label. Numerical features are standardised using z-score normalisation, and categorical features are transformed using one-hot encoding. The final encoded feature matrix is used consistently across xRFM, XGBoost, LightGBM, and MLP to avoid data-processing differences between models.

\subsubsection{Evaluation Metrics}

For regression tasks, we report Root Mean Squared Error (RMSE), where lower values indicate better predictive accuracy. For binary classification tasks, we report accuracy and AUC-ROC. AUC-ROC is especially important for the imbalanced classification datasets because it evaluates ranking ability across thresholds. We also record total training time and per-sample inference time to compare computational efficiency.

\subsubsection{xRFM Configuration}
\label{sec:xrfm-config}

The xRFM experiments use the Python \texttt{xrfm} package through the shared project training pipeline. The model is initialised as \texttt{xRFM(use\_diag=True, device=DEVICE)}, where \texttt{DEVICE} is CUDA if available and CPU otherwise. On GPU, the training script sets the number of fit iterations to $3$ and \texttt{M\_batch\_size=128} to control memory cost. Feature importance is extracted by reading the learned AGOP matrix or diagonal values from the trained xRFM model and ranking features by their AGOP-derived scores.

\begin{table}[H]
  \centering
  \caption{xRFM configuration used in the experiments.}
  \label{tab:xrfm-config}
  \begin{tabular}{ll}
    \toprule
    Component & Setting \\
    \midrule
    Model & \texttt{xRFM} \\
    AGOP mode & Diagonal AGOP, \texttt{use\_diag=True} \\
    Device & CUDA if available; otherwise CPU \\
    Random seed & 42 \\
    GPU fit iterations & 3 \\
    GPU AGOP batch size & 128 \\
    Data split & 60/20/20 train/validation/test \\
    Regression metric & RMSE \\
    Classification metrics & Accuracy and AUC-ROC \\
    Interpretability output & Top features ranked by AGOP diagonal score \\
    \bottomrule
  \end{tabular}
\end{table}

\subsubsection{XGBoost and LightGBM Configuration}
\label{sec:gbdt-config}

XGBoost and LightGBM are trained through the scikit-learn API using strong default settings from recent tabular benchmarking work~\citep{chen2016xgboost,ke2017lightgbm,holzmuller2024better}. Both use 1{,}000 maximum boosting rounds, learning rate $\eta=0.05$, row and column subsampling of 0.8, and early stopping with patience 20 on the validation set. XGBoost uses maximum tree depth 6, while LightGBM uses \texttt{num\_leaves=31}.

\subsubsection{MLP Configuration}
\label{sec:mlp-config}

The MLP baseline uses a two-hidden-layer network with hidden sizes $(128,64)$ and ReLU activation. It is trained with Adam, learning rate 0.001, batch size 256, maximum 500 iterations, and early stopping with patience 20. Input features are standardised before training. For classification tasks, predicted probabilities are used for AUC-ROC calculation.
